\section{Билет 7. Физический маятник. Уравнение малых колебаний и период колебаний физического маятника с выводом.}

\begin{definition}
\textbf{Физическим маятником} называется твёрдое тело, способное совершать колебания вокруг неподвижной горизонтальной оси, не проходящей через его центр масс. В качестве примера можно рассмотреть стержень, диск или любую другую конструкцию, закреплённую на горизонтальном шарнире.
\end{definition}

\begin{theorem}[Уравнение движения физического маятника]
Пусть маятник отклонён от положения равновесия на угол $\varphi$. Сила тяжести $m\vec{g}$, приложенная к центру масс $C$, создаёт вращательный момент относительно оси подвеса $O$:
\begin{equation}
    N = -mgl \sin\varphi, \label{eq:torque_phys}
\end{equation}
где $l$ --- расстояние от точки подвеса до центра масс. Знак «минус» указывает на то, что момент направлен в сторону, противоположную углу отклонения (возвращающий момент).

Запишем основное уравнение динамики вращательного движения:
\begin{equation}
    I \ddot{\varphi} = N,
\end{equation}
где $I$ --- момент инерции маятника относительно оси вращения. Подставляя \eqref{eq:torque_phys}, получаем точное нелинейное уравнение:
\begin{equation}
    I \ddot{\varphi} + mgl \sin\varphi = 0 \quad \Rightarrow \quad \ddot{\varphi} + \frac{mgl}{I} \sin\varphi = 0.
\end{equation}
Для малых углов отклонения ($\varphi \ll 1$ рад) справедливо приближение $\sin\varphi \approx \varphi$. Тогда уравнение принимает вид линейного однородного дифференциального уравнения второго порядка:
\begin{equation}
    \ddot{\varphi} + \frac{mgl}{I} \varphi = 0. \label{eq:small_phys}
\end{equation}
Вводя обозначение $\omega_0^2 = \frac{mgl}{I}$, приходим к каноническому уравнению гармонических колебаний:
\begin{equation}
    \ddot{\varphi} + \omega_0^2 \varphi = 0.
\end{equation}
\end{theorem}

\begin{property}
Период малых колебаний физического маятника определяется формулой:
\begin{equation}
    T = \frac{2\pi}{\omega_0} = 2\pi \sqrt{\frac{I}{mgl}}. \label{eq:period_phys}
\end{equation}
Согласно теореме Штейнера, момент инерции относительно оси подвеса равен $I = I_C + ml^2$, где $I_C$ --- момент инерции относительно оси, проходящей через центр масс параллельно оси вращения. Период зависит от распределения массы в теле (через $I_C$) и расстояния $l$ до центра масс.
\end{property}

\begin{remark}
\textbf{Приведённая длина.} Если ввести величину $l_{\text{пр}} = \frac{I}{ml}$, то формула периода \eqref{eq:period_phys} принимает вид $T = 2\pi \sqrt{l_{\text{пр}}/g}$, полностью аналогичный периоду математического маятника. Величину $l_{\text{пр}}$ называют \textbf{приведённой длиной физического маятника}. Математический маятник с длиной нити $l_{\text{пр}}$ будет колебаться с тем же периодом, что и данный физический маятник.
\end{remark}

\begin{figure}[htbp]
    \centering
    \begin{tikzpicture}[
        scale=1.2,
        force/.style={->, >=Stealth, thick},
        pivot/.style={fill=gray!50, draw=black, thick}
    ]
    % 1. Опора (стена/потолок)
    \fill[gray!30] (-2, 2.5) rectangle (2, 2.8);
    \draw[thick] (-2, 2.5) -- (2, 2.5);

    % Точка подвеса O
    \coordinate (O) at (0, 2.5);
    \draw[pivot] (O) circle (3pt);
    \node[above right] at (O) {$O$};

    % Параметры для расчетов
    \def\angle{30} % Угол отклонения
    \def\dist{3}   % Расстояние от оси до центра масс

    % 2. Вертикаль (пунктир)
    \draw[dashed, gray] (O) -- ++(0, -4.5);

    % 3. Тело маятника (вращается вокруг O)
    \begin{scope}[rotate=\angle]
        % Произвольная форма тела (пластина)
        \draw[fill=blue!10, thick, rounded corners] (-0.8, 0) rectangle (0.8, -4);

        % Ось симметрии тела
        \draw[gray, thin] (0, 0) -- (0, -4);

        % Центр масс C
        \coordinate (C) at (0, -\dist);
        \fill[black] (C) circle (2pt);
        \node[right=2pt] at (C) {$C$};

        % Расстояние a (или l)
        \draw[<->, gray] (0.5, 0) -- (0.5, -\dist) node[midway, right, gray, font=\small] {$a$};
    \end{scope}

    % 4. Силы и геометрия (рисуем вне scope, чтобы не вращались)

    % Координата C для рисования вектора силы тяжести (вычисляем вручную для точности)
    % C_x = a * sin(angle), C_y = 2.5 - a * cos(angle)
    \pgfmathsetmacro{\Cx}{\dist * sin(\angle)}
    \pgfmathsetmacro{\Cy}{2.5 - \dist * cos(\angle)}
    \coordinate (Cglobal) at (\Cx, \Cy);

    % Сила тяжести mg
    \draw[force, red] (Cglobal) -- ++(0, -1.5) node[midway, right, red] {$mg$};

    % Угол отклонения phi
    \draw[->, blue!70!black, thick] (0, 1.5) arc (90:90-\angle:1);
    \node[blue!70!black] at (0.6, 2) {$\varphi$};

    % Плечо силы тяжести (d = a * sin(phi)) - для вывода момента сил
    \draw[dashed, orange] (Cglobal) -- (0, \Cy) node[midway, above, orange, font=\small] {$d$};
    \draw[dashed, orange] (O) -- (0, \Cy);

    % Момент инерции J (или I)
    \node[align=left, font=\small, fill=white, opacity=0.8] at (-2.5, 0)
        {\textbf{Параметры:}\\
        $C$ --- центр масс\\
        $a$ --- расстояние $OC$\\
        $J$ --- момент инерции\\
        \textcolor{orange}{$d$} --- плечо силы};

    \end{tikzpicture}
    \caption{Физический маятник: твердое тело, совершающее колебания вокруг неподвижной горизонтальной оси $O$, не проходящей через центр масс $C$. Возвращающий момент сил $M = -mga\sin\varphi$.}
    \label{fig:physical_pendulum}
\end{figure}
